\section{Steady state}
\label{sec:steady}

\paragraph{References.} Table~\ref{tab:refs} lists the reference of each case,
how it was obtained, and the independent solutions that check it. The
primary reference is the candidate with the smallest own uncertainty:
\begin{itemize}
\item Langenbuch: the three-level limit of \ffn's degree-3 refinements (observed
  order 3.3), confirmed by degree 4 and 5 to \pending{0.34~pcm};
\item HP-MR: \ffn's finest run, where FE2 $r=2$ and FE3 $r=1$ agree to the
  printed digit; \garnet's mesh limit (order 2.2) is \pending{0.9~pcm} away;
\item IAEA-3D: see the paragraph on IAEA below.
\end{itemize}
Where \garnet's sequence is not yet asymptotic, its limit is listed but not
used. On Langenbuch it shows an order of 0.9 against a nominal 2: the error
there is dominated by the re-entrant corners of the staircase boundary. A
run that \emph{is} the reference is not a comparison row.

\begin{table}[h]
\tablefile{references}
\caption{References and their checks. Uncertainty: the largest of the
extrapolation step, the spread of the independent high-degree solutions,
and (for an asymptotic series) half the cross-code difference.}
\label{tab:refs}
\end{table}

\paragraph{Time to accuracy.} Table~\ref{tab:tta} gives, for each target, the
cheapest measured run of each platform. Figures~\ref{fig:iaea}--\ref{fig:hpmr}
show every run, with each platform's Pareto front, and
Tables~\ref{tab:runs-iaea}--\ref{tab:runs-hpmr} list them all.

\begin{table}[h]
\tablefile{tta_summary}
\caption{Time to accuracy (solve time, s): the cheapest run of each platform
that meets the target ($|\Delta k|$ and the largest node-power error; for the
transient, the largest power error over the run), and the ratio of \ffn's
better degree to \garnet. ``--'': no measured run meets the target.
\pending{The IAEA-3D rows use the provisional reference.}}
\label{tab:tta}
\end{table}

\paragraph{What the steady results show.}
\begin{itemize}
\item \textbf{Per core, \ffn's high-order elements win, often by an order of
  magnitude.} On the HP-MR at 10~pcm and 1\%, \ffn{} needs 89~s (FE3 $r=0$).
  \garnet{} on one core needs 2000~s ($r=6$, 10.5 million unknowns): 22 times
  slower. On Langenbuch the two are comparable at 10~pcm (45 against 66~s).
  \garnet's cost per unknown is lower (Section~\ref{sec:discussion}), but its
  second-order finite volumes need far more unknowns for the same error,
  above all for the node powers.
\item \textbf{\garnet's node powers converge slowly at the fuel--reflector
  corners.} On IAEA the largest node-power error falls roughly as
  \pending{$h^{1.6}$}, and at $r=16$ (37.5 million unknowns) it is still
  \pending{2.6\%}, while $k$ is within 10~pcm. The worst nodes are at the edge of the top fuel layer, where four
  of the eight surrounding nodes are radial or axial reflector, at a quarter
  of the average power. That is the corner of the fuel region, where a
  cell-centred scheme with harmonic-mean coefficients loses order. On
  Langenbuch the worst node, in the same position, is 0.56\% off at $r=16$.
\item \textbf{On the GH200, \garnet{} reaches the 10~pcm, 1\% target first}
  wherever it reaches it at all: 7.0~s on the HP-MR (13 times faster than \ffn{}
  on one core) and 0.35~s on Langenbuch (130 times). Its 1~pcm, 0.1\% target
  is out of reach of the meshes that fit (HP-MR $r=12$, 84 million unknowns,
  is 1.1~pcm and 0.16\% off). That target is close to the references' own
  resolution (0.3 to 0.4~pcm), so it separates the codes only coarsely.
\end{itemize}

\paragraph{IAEA-3D.} \pending{Until \ffn{} FE3 $r=2$ is in, the reference
is provisional: \garnet's GH200 mesh limit ($r$ = 8, 12, 16, observed order
1.57), uncertain to 3.2~pcm, with the node powers extrapolated at the same
order. \ffn's finest run, FE2 $r=2$, lands 0.6~pcm from it and within 0.09\%
in every node, so the two codes agree independently on $k$ and on the node
powers. With FE3 $r=2$, \ffn's degree-3 series becomes the reference, as on
Langenbuch.}
\begin{itemize}
\item \ffn{} meets 10~pcm and 1\% at FE3 $r=1$ (2.0 million unknowns,
  \pending{665~s}) and at FE2 $r=2$ (4.8 million, \pending{658~s}).
\item No \garnet{} run meets it, on either platform. At $r=16$ (37.5 million
  unknowns) $k$ is \pending{6.5~pcm} off, but the largest node-power error is
  still \pending{2.6\%}. The GH200 solves that mesh in 7.6~s, one core in
  2120~s. The small $\Delta k$ at $r=2$ ($-2.9$~pcm) is a cancellation of
  errors: its node powers are 31\% off.
\item \pending{The 1~pcm, 0.1\% target lies below the provisional
  reference's uncertainty, so the IAEA row for it in Table~\ref{tab:tta} is
  not yet meaningful.}
\end{itemize}

\begin{figure}[h]
\figfile{pareto_iaea3d}{0.95}
\caption{IAEA-3D (\pending{provisional reference}): $|\Delta k|$ and the largest node-power error against solve
time. Faint: every run; line: each platform's Pareto front; dashed: the
targets; grey: below the reference's uncertainty (hollow markers sit there).}
\label{fig:iaea}
\end{figure}
\begin{figure}[h]
\figfile{pareto_langenbuch3d}{0.95}
\caption{3-D Langenbuch at $t=0$, as Figure~\ref{fig:iaea}.}
\label{fig:lmw}
\end{figure}
\begin{figure}[h]
\figfile{pareto_hpmr3d}{0.95}
\caption{HP-MR, hex-homogenized, 11 groups, as Figure~\ref{fig:iaea}.}
\label{fig:hpmr}
\end{figure}
